\documentclass[11pt]{article}
\usepackage{fontspec}
\usepackage[french,english]{babel}
\usepackage{amsmath}
\usepackage{amssymb}
\usepackage{amsthm}
\usepackage{mathrsfs}
\usepackage{tikz-cd}
\usepackage{enumitem}

\newcommand{\printedpage}[1]{\marginpar{\footnotesize p.~#1}}
\newcommand{\msp}[1]{\marginpar{\footnotesize [#1]}}
\newcommand{\ux}{\underline{x}}
\newcommand{\cC}{\mathcal{C}}
\newcommand{\Ob}{\operatorname{Ob}}
\newcommand{\card}{\operatorname{card}}
\newcommand{\id}{\mathrm{id}}
\newcommand{\Hom}{\operatorname{Hom}}
\newcommand{\Ind}{\operatorname{Ind}}
\newcommand{\Img}{\operatorname{Im}}
\newcommand{\Ker}{\operatorname{Ker}}
\newcommand{\defeq}{\overset{\text{\upshape def}}{=}}

\begin{document}

\setcounter{footnote}{183}

\printedpage{169}
% head: A. Grothendieck : Dérivateurs, Ch. XVIII, Édition des Universités de Montpellier II et Paris VII
which is essentially surjective, but for lack of knowing that $\mathcal{E}(\pi) \subset \mathcal{E}_\pi$ in the cofibrant case, it is not clear that this functor is an equivalence of categories.

\textbf{4.15.7.} I must postpone the proof of lemma 4.15.3 (p. 428), assuming 4.8.13. (Not yet proved -- I am going to come back below to the proof of 4.8.13.) We are going to use the

\textbf{Lemma 4.15.7.1.} \textit{Let $\cC$ be a category, $\ux = (x_i)_{i\in I}$ an ind-object of $\cC$. For $\ux$ to be essentially constant, i.e. isomorphic, in} $\Ind(\cC)$\textit{, to the constant ind-object defined by an object}

\msp{page 441}
\textit{of $\cC$, it is necessary and sufficient that there exist a cofinal functor $\varphi : J \longrightarrow I$, with $J$ filtering, in such a way that the reindexed system $(x_j)_{j\in J}$, with $x_j = x_{\varphi(j)}$, have the following property: for every $j \in J$, one can find a projector $p_j$ in $x_j$, in such a way that:}
\begin{itemize}
\item[\textit{1°)}] \textit{For every arrow $\lambda : j \longrightarrow j'$ in $J$, we have commutativity in}
\[
\begin{tikzcd}
x_j \arrow[r, "x_\lambda"] \arrow[d, "p_j"] & x_{j'} \arrow[d, "p_{j'}"] \\
x_j \arrow[r, "x_\lambda"] & x_{j'}
\end{tikzcd}
\]
\textit{(where the $x_\lambda$ are the transition arrows), i.e. $p_{j'}x_\lambda = x_\lambda p_j$.}
\item[\textit{2°)}] \textit{For every $j \in \Ob J$, there exists an arrow $\lambda : j \longrightarrow j'$ in $J$ such that $p_{j'}x_\lambda = x_\lambda p_j = x_\lambda$.}
\item[\textit{3°)}] \textit{For every $j \in \Ob J$, the projector $p_j$ in $x_j$ admits an image $x_j^\natural$.} % [sic] un image -> une image, gender agreement
\item[\textit{4°)}] \textit{For every arrow $\lambda : j \longrightarrow j'$ in $J$, the arrow $x_\lambda^\natural : x_j^\natural \longrightarrow x_{j'}^\natural$ induced by $x_\lambda$ (thanks to 1°) is an isomorphism.}
\end{itemize}
\textit{Moreover, if we can suppose $J$ of the form $i_0\backslash I$, the cofinal functor $J \longrightarrow I$ being the canonical functor $i_0\backslash I \longrightarrow I$ ($i_0$ being a suitable fixed object of $I$)} \texttt{[incomplete sentence]}

\msp{page 442}
Let us note at once two corollaries.

\textbf{Corollary 4.15.7.2.} \textit{Suppose that $\cC$ is without projectors. Then $\ux = (x_i)_{i\in I}$ is an essentially constant ind-object if and only if there exists a cofinal functor, $J \longrightarrow I$, with $J$ filtering (\footnote{and we can suppose $J = i_0\backslash I \xrightarrow{\text{can.}} I$}), in such a way that the reindexed direct system $(x_j)_{j\in J}$ is isomorphic} (in the category $\underline{\mathrm{Hom}}(J,\cC)$) \textit{to a} constant \textit{direct system (i.e. defined by a constant functor $J \longrightarrow \cC$: all the $x_j$ equal to one and the same $x$, and the transition morphisms identities).}

Sufficiency is clear, for necessity we apply 4.15.7.1, which tells us that we can reindex so as to satisfy conditions 1°) to 4°) (for suitable projectors $p_j$ in the $x_j$). But the hypothesis on $\cC$ implies that $p_j = \id_{x_j}\ \forall\, j \in \Ob J$, and conditions 1°, 2°, 3° are satisfied tautologically for such a choice, condition 4° means that the transition morphisms are isomorphisms. But it is immediate that for a given direct system $(x_j)_{j\in J}$, with $J$ 1-\textit{connected} (e.g. $J$ filtering), the condition that the $x_j$ are isomorphisms % [sic] x_j -> x_\lambda, the transition morphisms are meant
is necessary and sufficient for this system to be isomorphic to a constant direct system, whence the conclusion.
% foot: Version 17 décembre 2025

\printedpage{170}
% head: A. Grothendieck : Dérivateurs, Ch. XVIII, Édition des Universités de Montpellier II et Paris VII
\msp{page 443}
\textbf{Corollary 4.15.7.3.} \textit{Suppose that $I$ is large relative to $\pi_0$, where $\pi_0$ is an infinite cardinal such that $\pi_0 \geqslant \card \cC$. Then $(x_i)_{i\in I}$ is essentially constant if and only if we can find a cofinal functor $J \longrightarrow I$ with $J$ large relative to $\pi_0$ (and ordered if we like) (\footnote{The proof (slightly rearranged) shows that we can take for $J$ a subcategory (not full) of a category $i_1\backslash I$, for a suitable object $i_1$ of $I$. Thus we have $\card J \leqslant \pi$ if $\card I \leqslant \pi$, $\pi$ being a given infinite cardinal.}), in such a way that the reindexed system $(x_j)_{j\in J}$ satisfies the following conditions:}
\begin{itemize}
\item[\textit{1°)}] \textit{The map $j \longmapsto x_j$ from $\Ob J$ to $\Ob\cC$ is constant, i.e. there exists an $x_0 \in \Ob\cC$ such that $x_j = x_0\ \forall\, j \in \Ob J$.}
\item[\textit{2°)}] \textit{There exists a projector $p$ in $x_0$ such that for $j < j'$ in $J$, we have $\tau_{j',j} = p$.}
\item[\textit{3°)}] \textit{The projector $p$ admits an image in $\cC$.}
\end{itemize}
\textsc{Proof}. These conditions obviously imply those of 4.15.7.1, taking $p_j = p\ \forall\, j \in \Ob J$, at least if $J$ has no greatest element (which causes a difficulty for 2°). But if $J$ has a final object, the same is true of $I$, and it is clear then that every direct system indexed by $I$ is essentially constant. So, in these cases, conditions 1°) 2°) 3°) above imply that $\ux$ is essentially constant. Converse\msp{page 444}ly, suppose that $\ux$ is essentially constant, and let us prove that we can reindex as indicated. The case where $I$ has a final object is trivial, we shall take $J = e$, $p = \id$. So suppose that $I$ has no final object. Since $I$ large relative to $\pi_0$ and $\pi_0 > \card\Ob\cC$, we see that there exists an $x_0 \in \Ob\cC$ such that the full subcategory $I' = I_{x_0}$ of $I$, consisting of the $i$ such that $x_i = x_0$, is cofinal in $I$. Since $(x_i)_{i\in I}$ is essentially constant, the same is true of $(x_{i'})_{i'\in I'}$, and we can apply criterion 4.15.7.1 to it, with \marginpar{$\Vert$}% double vertical bar printed in the left margin beside the two following lines
$J \longrightarrow I'$ cofinal. The proof will show moreover that if $I'$ is large relative to $\pi_0$ (an infinite cardinal with no other conditions), we can suppose $J$ large relative to $\pi_0$ (\footnote{take this into account in the proof of 4.15.7.1.}). Moreover, it is clear that we can suppose $J$ ordered. Since $\pi_0 \geqslant \card\Hom(x_0,x_0)$, we see again that there exists a projector $p$ in $x_0$ such that the subset $J_1 = J'$ of $J$, consisting of the $j$ such that $p_j = p$, is cofinal. Being cofinal, it too is large relative to $\pi_0$. So we can suppose that all the $p_j$ are equal to one and the same $p$. Conditions 1°-4° become:
\begin{itemize}
\item[1°)] $p\tau_{j'j} = \tau_{j'j}p$ for $j \leqslant j'$.
\item[2°)] For every $j$, exists $j' \geqslant j$ with
\[
p\tau_{j'j} \;=\; \tau_{j'j}p \;=\; \tau_{jj'}\ . % [sic] \tau_{jj'} -> \tau_{j'j}, index order used everywhere else
\]
\item[3°)] The image $x_0^\natural$ of the projector $p$ in $x_0$ exists.
\item[4°)] The endomorphism $\tau^\natural_{j'j}$ of $x_0^\natural$ is an isomorphism, for $j \leqslant j'$.
\end{itemize}
\msp{page 445}
Now we have the

\textbf{Lemma 4.15.7.4.} \textit{Let $\cC$ be a category, $x_0$ an object of $\cC$, $p$ a projector in $x_0$, $\tau$ an endomorphism of $x_0$ satisfying the relations}
\[
\tau p \;=\; p\tau \;=\; \tau\ , \tag*{$(*)$}
\]
% foot: Version 17 décembre 2025
\printedpage{171}
% head: A. Grothendieck : Dérivateurs, Ch. XVIII, Édition des Universités de Montpellier II et Paris VII
\textit{and such that the automorphism of $x_0^\natural \defeq$} $\Img p$ \textit{(defined e.g. as an object of $\cC^\wedge$, $x_0^\natural =$} $\Ker(\id_{x_0},p)$ % [sic] "(défini ... Ker(id_{x_0},p)" -> a closing parenthesis is missing, the parenthesis opened before "défini" is never closed
\textit{be the identity. Then we have $\tau = p$.}

\textsc{Proof of} 4.15.7.4. We are reduced at once, by the standard procedure, to the case where $\cC = \mathrm{Ens}$, where the verification is immediate: the automorphisms $\tau$ of the set $E$ $(= x)$ which satisfies $(*)$ correspond one-to-one to the endomorphisms $u$ of $E^\natural = pE$, via
\[
\begin{cases}
u \;\longmapsto\; \sigma u \rho \\
\tau \;\longmapsto\; \tau|E^\natural
\end{cases}
\qquad
\text{(where }
\begin{tikzcd}[column sep=small]
E^\natural \arrow[r, hook, shift left, "\sigma"] & E \arrow[l, shift left, "\rho"]
\end{tikzcd}
\text{ are the canonical arrows)}\ ,
\]
so $\tau|E^\natural = \id \Longrightarrow u = \sigma\rho = p$.

Let us choose a $j_0 \in J$, and consider the homomorphisms
\[
\tau^\natural_{j,j_0} \;:\; x_0^\natural \longrightarrow x_0^\natural
\]
induced by the $\tau_{j,j_0}$ with $j \geqslant j_0$. Here again, since the set of possible isomorphisms is of

\msp{page 446}
cardinality $\leqslant \pi_1$, and $J/j_0$ is large relative to $\pi_0$, we see that there exists $J' \subset j_0\backslash J$, cofinal in $j_0\backslash J$, hence in $J$, such that for $j \in J'$, the $\tau^\natural_{j,j_0}$ are equal to one and the same $\tau : x_0^\natural \xrightarrow{\sim} x_0^\natural$. So if $j', j'' \in J'$, with $j' \leqslant j''$, we have
\[
\underbrace{\tau_{j''j_0}}_{\tau} \;=\; \tau_{j''j'}\,\underbrace{\tau_{j'j_0}}_{\tau}\ ,
\]
whence
\[
\tau_{j''j'} \;=\; \id_{x_0^\natural}\ ,
\]
since the arrow $\tau$ is an isomorphism. So the direct system reindexed by $J'$, condition 4° is replaced by the stricter condition \fbox{4° bis) $\tau^\natural_{j',j} = \id_{x_0^\natural}$.} But condition 2° then becomes, thanks to 4.15.7.4,
\[
\boxed{\text{2° bis) For every } j\text{, exists } j' \geqslant j \text{ with } \tau_{j',j} = p}\ .
\]
So, we can suppose 1°, 2° bis, 3°, 4° bis (we again denote by $J$, instead of $J'$, the indexing ordered set large relative to $\pi_0$).

Let $K$ be the ordered set defined by

\msp{page 447}
\[
\begin{cases}
\Ob K \;=\; \Ob J \\
\text{if } j,j' \in \Ob J\text{, we have } j \overset{K}{\leqslant} j' \text{ if and only if } j \overset{\mathrm{J}}{=} j'\text{, or } j \overset{J}{<} j' \text{ and } \tau_{j',j} = p\,.
\end{cases}
\]
It is immediate that this is an \textit{order relation} on $J$ (since $p^2 = p$), finer than the initial order relation. The direct system induced on $K$ satisfies conditions 1°, 2°, 3° of 4.15.7.3. On the other hand, $K \longrightarrow J$ is cofinal, hence so is the composite $K \longrightarrow J \longrightarrow I$. It remains to prove that $K$ is large relative to $\pi_0$, hence that every subset $K' \subset K$ such that $\card K' \leqslant \pi_0$ has an upper bound in $K$. It has an upper bound \textit{in} $J$, say $j$. By hypothesis 2 bis, there exists $j' \geqslant j$ such that $\tau_{j'j} = p$, and we can suppose $j' > j$. (Indeed, if $\tau_{j'j} = p$, we see that
% foot: Version 17 décembre 2025
\printedpage{172}
% head: A. Grothendieck : Dérivateurs, Ch. XVIII, Édition des Universités de Montpellier II et Paris VII
\msp{page 448}
the initial form 2°) of condition 2° bis that we have $\tau_{j''j} = p$ for every $j'' \geqslant j'$, so if $j$ is not a final object of $J$, % [sic] par -> pas, misprint
we can suppose $j' > j$. But $J$ has no final object, for otherwise $I$ would have one, a case we had eliminated. So there exists $j'$ such that $j' \overset{K}{>} j$.) Now we have this (which we apply to the case $i \in K'$):
\[
i \leqslant j \ \text{ and } \ j \overset{K}{<} j' \qquad \Longrightarrow \qquad i \overset{K}{<} j'\ .
\]
\[
\underbrace{x_i}_{= x_0} \xrightarrow{\tau_{j,i}} x_j \xrightarrow{\tau_{j',j}=p} x_{j'}\ ,
\]
\[
\tau_{j'i} \;=\; \tau_{j'j}\tau_{ji} = \underbrace{p\tau_{ji}}_{\overset{\text{1°}}{=}\,\tau_{ji}p} \;\overset{?}{=}\; p\ ,
\]
we have indeed
\[
\left.
\begin{aligned}
\tau_{j',i} &\;=\; p\tau_{j',i} \\
\tau_{j',i} &\;-\; \tau_{j',i}p % [sic] - -> =, the second relation of (*) is meant
\end{aligned}
\right\}
\quad \overset{\text{by 4.15.7.4}}{\Longrightarrow} \quad % Doubt resolved against PDF leaf 14: no underline; the two strokes belong to the implication arrow.
\tau_{j',i} \;=\; p\ .
\]
This completes the proof of corollary 4.15.7.3.

\textbf{Corollary 4.15.7.5} (of 4.15.7.3). \textit{Let $(x_i)_{i\in I}$ be an essentially constant direct system large relative to $\pi_0 \geqslant$} $\card\cC$\textit{, in $\cC$} without projectors. \textit{Then there exists $J \longrightarrow I$ cofinal, with $J$ ordered large relative to $\pi_0$, and such that the direct system induced on $J$ is} constant.

\textsc{Let us prove lemma} 4.15.7.1. We must make explicit the existence of two arrows in $\Ind(\cC)$,
\[
z \underset{\beta}{\overset{\alpha}{\rightleftarrows}} \text{"}\varinjlim_{i\in I}\text{"} x_i \;=\; \ux\ ,
\]
with $\beta\alpha = \id_z$, $\alpha\beta = \id_{\ux}$. The datum

\msp{page 449}
of $x$ is equivalent to that of an element of $\varinjlim\limits_{i\in I} \Hom(z,x_i)$, and is therefore given by an % [sic] x -> \alpha, parallel with "la donnée de \beta" below
\[
\alpha_{i_0} \;:\; z \longrightarrow x_{i_0}\ ,
\]
the datum of $\beta$ is equivalent to that of an element of $\varprojlim\limits_{i\in I} \Hom(x_i,z)$, i.e. of arrows
\[
\beta_i \;:\; x_i \longrightarrow z\ ,
\]
with the usual compatibility condition. At the cost of replacing $I$ by $I'$, with $I' \longrightarrow I$ cofinal, we can suppose $I$ ordered (and large relative to $\pi_0$, if the initial $I$ was), and the compatibility condition on the $\beta_i$ reads
\[
\beta_i\tau_{ji} \;=\; \beta_j \qquad \text{if } i \leqslant j \text{ in } I,
\]
where $\tau_{ji} : x_i \longrightarrow x_j$ denotes the transition arrow. We must make explicit, in terms of these data, the conditions
\[
\alpha\beta \;=\; \id_z\ , \qquad \beta\alpha \;=\; \id_{\ux}\ .
\]
% foot: Version 17 décembre 2025

\printedpage{173}
% head: A. Grothendieck : Dérivateurs, Ch. XVIII, Édition des Universités de Montpellier II et Paris VII
\[
\begin{tikzcd}
z \arrow[r, "\alpha_{i_0}"] \arrow[dr] & x_{i_0} \arrow[r, "\tau_{i,i_0}"] \arrow[d, "\beta_{i_0}"] & x_i \arrow[dl, "\beta_i"] \\
 & z &
\end{tikzcd}
\]
The relation $\alpha\beta = \id_z$ is equivalent to
\[
\beta_{i_0}\alpha_{i_0} \;=\; \id_z\ , \tag*{$(*_{i_0})$}
\]
it implies moreover, on setting for $i \geqslant i_0$,
\[
\alpha_i \;\defeq\; \tau_{i,i_0}\alpha_{i_0} \;:\; z \longrightarrow x_i\ ,
\]
that we have for $i \geqslant i_0$
\[
\beta_i\alpha_i \;=\; \id_z \qquad \forall\, i \geqslant i_0\ . \tag*{$(*_i)$}
\]

\msp{page 450}
Let us then set
\[
p_i \;=\; \alpha_i\beta_i \;:\; x_i \longrightarrow x_i\ ,
\]
it is clear that we have
\[
\begin{cases}
p_i^2 \;=\; p_i \qquad \text{($p_i$ a projector)} \\
\alpha_i \;:\; z \longrightarrow x_i \qquad \text{induces an isomorphism} \\
\qquad z \xrightarrow{\sim} x_i^\natural \;\defeq\; \Img p_i \;=\; \Ker(\id_{x_i},p_i)\ .
\end{cases}
\]
The desired cofinal ordered set $J$ is $i_0\backslash I$, it is filtering, and is large relative to $\pi_0$ if $I$ is. To tell the truth, the reduction to the case $I$ ordered is unnecessary, we can take $J = i_0\backslash I$ directly, it is immediate that $J$ is filtering ($I$ being so), and $i_0\backslash I \longrightarrow I$ cofinal, finally that if $I$ is large relative to an infinite cardinal $\pi_0$, the same is true of $i_0\backslash I$.

Let us express that $\beta\alpha = \id_{\ux}$; now we see that $\beta\alpha$ is defined by the endomorphism $(p_j)_{j\in J}$ of the direct system $(x_j)_{j\in J}$ -- the fact that this really is an endomorphism of this system, i.e. that we have relation 1°) of 4.15.7.1,
\[
p_{j'}x_\lambda \;=\; x_\lambda p_\lambda \qquad \text{if } \lambda : j \longrightarrow j' \text{ in } J, % [sic] p_\lambda -> p_j, cf. relation 1°) of 4.15.7.1
\]
is immediate, it reads
\[
\alpha_{j'}\underbrace{\beta_{j'}x_\lambda}_{\beta_j} \;=\; \underbrace{x_\lambda\alpha_j}_{\alpha_{j'}}\beta_j\ , \qquad \text{now both sides are equal to } \alpha_{j'}\beta_j.
\]
% foot: Version 17 décembre 2025

\end{document}
